% ------------------------------------------------------------
% Page 112
% ------------------------------------------------------------

« Le château » : tient une place à part dans l’histoire de Ferrussac.

A l’origine il faisait sans doute partie du domaine de Roquedols. C’est en 1842 que
\textbf{Jacques Cambessédès}, botaniste connu en devint propriétaire.

Né à Montpellier d’un père originaire du Vigan et de famille protestante, au 17\textsuperscript{e} siècle, puis « nouveau
converti », il commença ses études chez les Oratoriens à Tournon puis fit sa médecine à Montpellier,
tout en poursuivant des études de botaniste.
Parti à Paris, militaire, il devint aide de camp de son oncle le général Mathieu Dumas. Lors de la
Révolution de 1830 (Les Trois Glorieuses) il défend l’ordre en face des révolutionnaires. C’était, dit-
on un brillant causeur, beau cavalier et ardent chasseur.
Il fit plusieurs randonnées dans les Baléares et les Pyrénées et les via avec les plus célèbres botanistes
de l’époque, Jussieu, St Hilaire, Jacquemont, Planchon de Montpellier…

A la mort de son père en 1835 il hérita de la vaste propriété de \textbf{Pradines}, sur le Causse Noir, « à
quelques pas du hameau de Lanuéjols et non loin de la petite ville de Meyrueis, écrit Emile Planchon,
dans une longue notice nécrologique (1), Il en fit une exploitation modèle, avec assolements,
amélioration du cheptel ovin par l’achat de béliers mérinos, importation de chevaux et mulets, achat
d’une batteuse anglaise et amélioration des charrues… »
Il était aussi lieutenant de Louveterie, « les bois des Cévennes retentirent du bruit de son cor et des
aboiements de sa meute de chiens… ».

Il avait été élu maire de Lanuéjols puis Conseiller Général lors de 4 scrutins.

Il trouva, près de Lanuéjols, des filons de charbon qui furent exploités dans la 2\textsuperscript{o} moitié du 19\textsuperscript{e} siècle,
mais découvrit aussi et décrivit l’adonis vernal ou différentes spirées du Causse, insista sur
l’importance des pâturages calcaires pour la croissance du cheval.
« Accessible à tous, il donnait à tous ses conseils et était adoré de ses serviteurs…. Le
désintéressement était un des traits de sa nature… L’hiver, pour un montagnard du Causse, est une
rude étape à franchir : la chasse en dissipait les ennuis. Avec l’été, venaient les travaux de plein air ;
l’automne, saison sociale par excellence, amenait à Pradines et plus tard à Ferrussac des hôtes
aimables, et la modeste demeure du Causse ou du vallon devenait pour quelques mois, le rendez-vous
d’une société choisie »

C’est en 1842 qu’il acheta le « Château » où il ne vint résider que vers 1855, semble-t-il, voulant ainsi
épargner à sa femme les trop rudes hivers du Causse et rendre la solitude moins pesante à Ferrussac,
avec la proximité de la ville de Meyrueis et du château de Roquedols.
Il fut élu conseiller municipal et nommé Maire de Meyrueis en 1860.

\footnote{Jules Emile Planchon : Médecin-Pharmacien, Directeur du Jardin des Plantes de Montpellier il identifia le
phylloxéra et sauva le vignoble français en sélectionnant des variétés de vignes résistantes. Il découvrit et
décrivit le kiwi !}
